\subsection{Что такое Python и зачем он нужен?}

\textbf{Мотивирующий вопрос:} Как объяснить компьютеру или роботу, что именно нужно сделать, используя понятный и лаконичный язык?

Python — это высокоуровневый язык программирования, который был создан Гвидо ван Россумом и впервые выпущен в 1991 году. Назван он не в честь змеи, а в честь популярного британского комедийного шоу <<Летающий цирк Монти Пайтона>>.

\subsubsection*{Почему Python?}
\begin{itemize}
    \item \textbf{Простота:} Синтаксис Python похож на обычный английский язык.
    \item \textbf{Универсальность:} Используется в веб-разработке, науке, AI и робототехнике.
    \item \textbf{Мощь:} Огромное количество библиотек (OpenCV, NumPy, TensorFlow) для решения сложных задач.
\end{itemize}

Для программирования квадрокоптеров Pioneer мы выбрали Python, так как он позволяет легко взаимодействовать с оборудованием и обрабатывать видеопоток с камеры.

\subsubsection*{Первая программа}
Традиционно изучение любого языка начинается с программы, которая выводит приветствие на экран.

\begin{codefile}[python]{examples/python/01_hello_world.py}
\end{codefile}

Здесь функция \texttt{print()} выводит текст, заключенный в кавычки, в консоль.

\begin{taskbox}[Привет, мир!]
Напишите программу, которая выводит ваше имя и фамилию на экран.
\end{taskbox}
